\noindent\markforreview
\vspace{1em}

\begin{center}
    \includegraphics[width=0.42\linewidth]{hard_10_f8.png}
\end{center}

In the figure, parallel lines \(a\) and \(b\) are intersected by lines \(c\), \(d\), and \(e\).

If \(z = 49^\circ\), \(y = 136^\circ\), and \(v < z\), which statement about \(x\) and \(w\) must be true?

\vspace{0.75em}

\choicebox{A}{\(x < w\)}
\choicebox{B}{\(x > w\)}
\choicebox{C}{\(x = w\)}
\choicebox{D}{\(x + w = 90\)}

\vspace{2em}
\newpage

\noindent\markforreview
\vspace{1em}

The functions \(f\) and \(g\) are defined by the given equations below, where \(x \geq 0\).

Which of the following equations displays, as a constant or coefficient, the minimum value of the function it defines?

\begin{itemize}
\item[I.] \(f(x) = 16(1.3)^{x+5}\)
\item[II.] \(g(x) = 16(1.09)(1.3)^{x+3}\)
\end{itemize}

\vspace{0.75em}

\choicebox{A}{I only}
\choicebox{B}{II only}
\choicebox{C}{I and II}
\choicebox{D}{Neither I nor II}

\vspace{2em}
\newpage

\noindent\markforreview
\vspace{1em}

The table shows values of \(x\) and corresponding values of \(g(x)\), where:
\[
g(x) = \frac{f(x)}{x + 3}
\]
and \(f(x)\) is a linear function.

What is the \(y\)-intercept of the graph of \(y = f(x)\) in the \(xy\)-plane?

\[
\begin{array}{c|c}
x & g(x) \\
\hline
-15 & 3 \\
-6 & 0 \\
9 & 5 \\
\end{array}
\]

\vspace{0.75em}

\choicebox{A}{\((0, -6)\)}
\choicebox{B}{\((0, 4)\)}
\choicebox{C}{\((0, 8)\)}
\choicebox{D}{\((0, 24)\)}

\vspace{2em}
\newpage

\noindent\markforreview
\vspace{1em}

The positive number \(a\) is 2,106\% of the sum of the positive numbers \(b\) and \(c\), and \(b\) is 81\% of \(c\).

What percent of \(b\) is \(a\)?

\vspace{0.75em}

\choicebox{A}{4,706\%}
\choicebox{B}{2,600\%}
\choicebox{C}{38.12\%}
\choicebox{D}{21.87\%}

\vspace{2em}
\newpage

\noindent\markforreview
\vspace{1em}

In the given equation, \(u\) is a positive constant.

The sum of the solutions to the equation is 1. What is the value of \(u\)?
\[
7x^2(3x - 11)(3x - u) = 0
\]

\vspace{0.75em}

\choicebox{A}{1}
\choicebox{B}{-8}
\choicebox{C}{\(\frac{11}{3}\)}
\choicebox{D}{\(\frac{8}{3}\)}

\vspace{2em}
\newpage

\noindent\markforreview
\vspace{1em}

Which of the following values of \(x\) satisfy the equation
\[
-3x^2 + 12x + 4 = 0
\]

\vspace{0.75em}

\choicebox{A}{0}
\choicebox{B}{\(-2\) and 2}
\choicebox{C}{\(2 - \frac{2\sqrt{6}}{3}\) and \(2 + \frac{2\sqrt{6}}{3}\)}
\choicebox{D}{\(2 - \frac{4\sqrt{3}}{3}\) and \(2 + \frac{4\sqrt{3}}{3}\)}

\vspace{2em}
\newpage

\noindent\markforreview
\vspace{1em}

A line intersects two parallel lines, forming four acute angles and four obtuse angles.

The measure of one of these eight angles is \((7x - 250)^\circ\).

The sum of the measures of four of the eight angles is \(k\).

Which of the following could **not** be equivalent to \(k\) for all values of \(x\)?

\vspace{0.75em}

\choicebox{A}{\(-14x + 1540\)}
\choicebox{B}{\(14x - 320\)}
\choicebox{C}{\(-28x + 1720\)}
\choicebox{D}{360}

\vspace{2em}
\newpage

\noindent\markforreview
\vspace{1em}

The triangle inequality theorem states that the sum of any two sides of a triangle must be greater than the length of the third side.

If a triangle has side lengths of 7 and 10, which inequality represents the possible lengths, \(x\), of the third side of the triangle?

\vspace{0.75em}

\choicebox{A}{\(x < 17\)}
\choicebox{B}{\(x > 17\)}
\choicebox{C}{\(3 < x < 17\)}
\choicebox{D}{\(x < 3\) or \(x > 17\)}

\vspace{2em}
